\chapter{Omtrek, oppervlakte, volume}\label{ch:g6:measure}

Hoe lang is het hek, hoe groot is het veld, hoeveel water gaat er in de tank?
Drie verschillende vragen, drie verschillende grootheden --- \omterm{def:g6:measure:perimeter}{omtrek}, \omterm{def:g6:measure:area}{oppervlakte}
en \omterm{def:g6:measure:volume}{volume} --- elk met haar eigen eenheden. Ze door elkaar halen is de meest
gemaakte fout in de meetkunde; dit hoofdstuk zet ze voor eens en altijd op hun
plaats.

\section{Lengtes en omtrek}

\begin{definition}[Omtrek]\label{def:g6:measure:perimeter}
De \emph{omtrek}\index{omtrek} van een figuur is de totale lengte van haar
rand. Je meet hem in lengte-eenheden: millimeter (mm), centimeter (cm), meter
(m), kilometer (km), met
\[
1 \text{ km} = 1000 \text{ m}, \qquad
1 \text{ m} = 100 \text{ cm}, \qquad
1 \text{ cm} = 10 \text{ mm}.
\]
\end{definition}

\begin{example}\label{ex:g6:measure:perimeter}
Een rechthoek met lengte $l$ en breedte $b$ heeft \omterm{def:g6:measure:perimeter}{omtrek}
\[
P = l + b + l + b = 2 \times (l + b).
\]
Voor een rechthoek van $7$ cm bij $4$ cm: $P = 2 \times (7 + 4) = 2 \times 11 =
22$ cm. Een vierkant met zijde $z$ heeft \omterm{def:g6:measure:perimeter}{omtrek} $4z$.
\end{example}

\begin{proposition}[Omtrek van een cirkel]\label{prop:g6:measure:circle}
De \omterm{def:g6:measure:perimeter}{omtrek} van een cirkel met \omterm{def:g3:shapes:shapes}{radius} $r$ is
\[
P = 2 \pi r ,
\]
waarbij $\pi \approx 3.14$ voor elke cirkel hetzelfde getal is.
\end{proposition}
\admitted

\begin{example}
Een ronde vijver heeft \omterm{def:g3:shapes:shapes}{radius} $5$ m. Zijn rand meet
$2 \times \pi \times 5 = 10\pi \approx 31.4$ m. Houd de precieze waarde
$10\pi$ zo lang mogelijk vast; rond pas aan het einde af.
\end{example}

\section{Oppervlaktes}

\begin{definition}[Oppervlakte]\label{def:g6:measure:area}
De \emph{oppervlakte}\index{oppervlakte} van een figuur meet het oppervlak dat
ze bedekt: hoeveel eenheidsvierkantjes er in passen. Eenheden: cm$^2$ (een
vierkantje met zijde $1$ cm), m$^2$, km$^2$, \dots\ Let op:
\[
1 \text{ m}^2 = 100 \times 100 \text{ cm}^2 = 10\,000 \text{ cm}^2
\]
--- een vierkante meter is een vierkant van $100$ cm bij $100$ cm, dus is elke
stap op de maatladder $100$ waard en niet $10$.
\end{definition}

\begin{omfigure}
\begin{tikzpicture}[scale=0.68]
  \draw[gray!50, thin] (0,0) grid (7,4);
  \draw[very thick, omDef] (0,0) rectangle (7,4);
  \fill[omThm!25] (0,3) rectangle (1,4);
  \draw[omThm, thick] (0,3) rectangle (1,4);
  \node[font=\small, omThm] at (2.6,3.5) {$1$ eenheidsvierkantje};
  \node[font=\small] at (3.5,-0.5) {$7 \times 4 = 28$ eenheidsvierkantjes};
\end{tikzpicture}

{\small De \omterm{def:g6:measure:area}{oppervlakte} van een rechthoek: $7$ kolommen van elk $4$
eenheidsvierkantjes, dus $7 \times 4 = 28$ vierkantjes in totaal. Daarom is
\omterm{def:g6:measure:area}{oppervlakte} $=$ lengte $\times$ breedte.}
\end{omfigure}

\begin{proposition}[De basisformules voor oppervlaktes]\label{prop:g6:measure:formulas}
\begin{center}
\renewcommand{\arraystretch}{1.5}
\begin{tabular}{l|l}
figuur & \omterm{def:g6:measure:area}{oppervlakte} \\
\hline
rechthoek ($l \times b$) & $A = l \times b$ \\
vierkant (zijde $z$) & $A = z^2$ \\[4pt]
\omterm{def:g6:shapes:triangles}{rechthoekige driehoek} (rechthoekszijden $a$, $b$) &
$A = \dfrac{a \times b}{2}$ \\[4pt]
schijf (\omterm{def:g3:shapes:shapes}{radius} $r$) & $A = \pi r^2$ \\
\end{tabular}
\end{center}
\end{proposition}

\begin{proof}[Bewijs voor de rechthoekige driehoek]
Twee kopieën van een \omterm{def:g6:shapes:triangles}{rechthoekige driehoek} met rechthoekszijden $a$ en $b$,
langs de schuine zijde aan elkaar geplakt, vormen een rechthoek $a \times b$. De
\omterm{def:g6:measure:area}{oppervlakte} van de driehoek is dus de \omterm{def:g3:division:half}{helft} van die van de rechthoek:
$\frac{ab}{2}$. (De formule voor de rechthoek is het tellen van vierkantjes
hierboven; de formule voor de schijf nemen we zonder bewijs aan, zie
\cref{ch:g7:areas}.)
\end{proof}

\begin{omfigure}
\begin{tikzpicture}[scale=0.75]
  \draw[thick, fill=omDef!15] (0,0) -- (4,0) -- (0,2.5) -- cycle;
  \draw[thick, omExo, dashed] (4,0) -- (4,2.5) -- (0,2.5);
  \draw[thin] (0.3,0) -- (0.3,0.3) -- (0,0.3);
  \node[below, font=\small] at (2,-0.1) {$a$};
  \node[left, font=\small] at (0,1.25) {$b$};
  \node[font=\small, omExo] at (2.8,1.7) {tweede kopie};
\end{tikzpicture}

{\small Een \omterm{def:g6:shapes:triangles}{rechthoekige driehoek} is de \omterm{def:g3:division:half}{helft} van een rechthoek: vandaar de
formule $\frac{a\times b}{2}$.}
\end{omfigure}

\begin{example}[Samengestelde figuren]\label{ex:g6:measure:composite}
Een L-vormige kamer is een rechthoek van $6$ m $\times$ $4$ m met een vierkante
hoek van $2$ m $\times$ $2$ m eruit. Haar \omterm{def:g6:measure:area}{oppervlakte}, stap voor stap:
\begin{enumerate}
  \item de volle rechthoek: $6 \times 4 = 24$ m$^2$;
  \item het weggehaalde vierkant: $2 \times 2 = 4$ m$^2$;
  \item wat overblijft: $24 - 4 = 20$ m$^2$.
\end{enumerate}
Haar \emph{\omterm{def:g6:measure:perimeter}{omtrek}} is geen $20$ van wat dan ook: als je om de L heen loopt, meet
de rand nog steeds $6 + 4 + 6 + 4 = 20$ m --- de twee sneden in de hoek
vervangen twee even lange stukken muur. Bij toeval hetzelfde getal, maar
vierkante meters voor het ene en meters voor het andere!
\end{example}

\begin{remark}[Dezelfde omtrek, andere oppervlaktes]
Twee figuren kunnen dezelfde \omterm{def:g6:measure:perimeter}{omtrek} en heel verschillende \omterm{def:g6:measure:area}{oppervlaktes} hebben:
een vierkant van $5 \times 5$ en een rechthoek van $9 \times 1$ hebben beide
\omterm{def:g6:measure:perimeter}{omtrek} $20$, maar \omterm{def:g6:measure:area}{oppervlaktes} $25$ en $9$. \omterm{def:g6:measure:perimeter}{Omtrek} en \omterm{def:g6:measure:area}{oppervlakte} zijn echt
onafhankelijke grootheden.
\end{remark}

\section{Volumes}

\begin{definition}[Volume]\label{def:g6:measure:volume}
Het \emph{volume}\index{volume} van een \omterm{def:g5:solids:def}{lichaam} meet de ruimte die het vult:
hoeveel eenheidskubusjes er in passen. Eenheden: cm$^3$, m$^3$, \dots\ met
$1$ m$^3 = 1\,000\,000$ cm$^3$ (elke stap op de ladder is $1000$ waard). Voor
vloeistoffen gebruikt men ook de \emph{liter}:
\[
1 \text{ L} = 1 \text{ dm}^3 = 1000 \text{ cm}^3,
\qquad
1 \text{ m}^3 = 1000 \text{ L}.
\]
\end{definition}

\begin{proposition}[Volume van een balk]\label{prop:g6:measure:box}
Een \omterm{def:g5:solids:def}{balk} met lengte $l$, breedte $b$ en hoogte $h$ heeft \omterm{def:g6:measure:volume}{volume}
\[
V = l \times b \times h .
\]
\end{proposition}

\begin{proof}
De onderste laag bevat $l \times b$ eenheidskubusjes (het plaatje bij
\cref{def:g6:measure:area}, met kubusjes), en er zijn $h$ van die lagen.
\end{proof}

\begin{omfigure}
\begin{tikzpicture}[scale=0.6]
  \foreach \y in {0,1} {
    \foreach \x in {0,...,3} {
      \foreach \z in {0,1} {
        \begin{scope}[shift={({\x+0.45*\z},{\y+0.3*\z})}]
          \draw[fill=omDef!12] (0,0) rectangle (1,1);
          \draw[fill=omDef!20] (0,1) -- (0.45,1.3) -- (1.45,1.3) -- (1,1)
            -- cycle;
          \draw[fill=omDef!28] (1,0) -- (1.45,0.3) -- (1.45,1.3) -- (1,1)
            -- cycle;
        \end{scope}
      }
    }
  }
  \node[below, font=\small] at (2.5,-0.3)
    {$4 \times 2 \times 2 = 16$ eenheidskubusjes};
\end{tikzpicture}

{\small Een \omterm{def:g5:solids:def}{balk} gevuld met eenheidskubusjes: $4 \times 2$ kubusjes per laag,
twee lagen.}
\end{omfigure}

\begin{example}\label{ex:g6:measure:aquarium}
Een aquarium meet $60$ cm bij $30$ cm bij $40$ cm (hoogte). \omterm{def:g6:measure:volume}{Volume}:
\[
V = 60 \times 30 \times 40 = 72\,000 \text{ cm}^3 = 72 \text{ dm}^3
= 72 \text{ L}.
\]
Stap voor stap voor de omzetting: $1000$ cm$^3$ maken één liter, en
$72\,000 \div 1000 = 72$.
\end{example}

\section{Oefeningen}

\begin{exercise}[$\star$]\label{exo:g6:measure:1}
Zet om: $3.5$ m in cm; $420$ mm in cm; $0.8$ km in m; $25\,000$ m in km.
\end{exercise}

\begin{exercise}[$\star$]\label{exo:g6:measure:2}
Reken de \omterm{def:g6:measure:perimeter}{omtrek} uit van: een rechthoek van $8$ cm $\times$ $3.5$ cm; een
vierkant met zijde $6.2$ cm; een driehoek met zijden $5$ cm, $7$ cm en $9$ cm.
\end{exercise}

\begin{exercise}[$\star$]\label{exo:g6:measure:3}
Een ronde atletiekbaan heeft \omterm{def:g3:shapes:shapes}{radius} $50$ m. Hoe lang is één rondje (precieze
waarde met $\pi$, en daarna afgerond op de meter)? Hoeveel rondjes maken minstens
$2$ km?
\end{exercise}

\begin{exercise}[$\star$]\label{exo:g6:measure:4}
Reken de \omterm{def:g6:measure:area}{oppervlakte} uit van: een rechthoek van $9$ cm $\times$ $4$ cm; een
vierkant met zijde $7$ m; een \omterm{def:g6:shapes:triangles}{rechthoekige driehoek} met rechthoekszijden $6$ cm
en $10$ cm; een schijf met \omterm{def:g3:shapes:shapes}{radius} $3$ cm (precieze waarde, en daarna afgerond op
de cm$^2$).
\end{exercise}

\begin{exercise}[$\star$]\label{exo:g6:measure:5}
Zet om: $3$ m$^2$ in cm$^2$; $45\,000$ cm$^2$ in m$^2$; $2.5$ L in cm$^3$;
$4\,500$ L in m$^3$.
\end{exercise}

\begin{exercise}[$\star$]\label{exo:g6:measure:6}
Een \omterm{def:g6:shapes:triangles}{rechthoekig} veld is $120$ m lang en $85$ m breed. Hoeveel meter hek is er
nodig om het te omheinen? En wat is zijn \omterm{def:g6:measure:area}{oppervlakte}?
\end{exercise}

\begin{exercise}[$\star$]\label{exo:g6:measure:7}
Reken het \omterm{def:g6:measure:volume}{volume} uit van een \omterm{def:g5:solids:def}{balk} van $5$ cm $\times$ $4$ cm $\times$ $10$ cm,
en van een \omterm{def:g5:solids:def}{kubus} met \omterm{def:g5:solids:def}{ribbe} $3$ cm.
\end{exercise}

\begin{exercise}[$\star$]\label{exo:g6:measure:8}
Teken twee verschillende rechthoeken met \omterm{def:g6:measure:perimeter}{omtrek} $16$ cm en reken hun
\omterm{def:g6:measure:area}{oppervlaktes} uit. Welke van jouw rechthoeken heeft de grootste \omterm{def:g6:measure:area}{oppervlakte}?
\end{exercise}

\begin{exercise}[$\star\star$]\label{exo:g6:measure:9}
Een T-vormige figuur bestaat uit een horizontale rechthoek van $10 \times 2$
boven een verticale van $2 \times 6$ (maten in cm). Reken haar \omterm{def:g6:measure:area}{oppervlakte} uit,
en daarna haar \omterm{def:g6:measure:perimeter}{omtrek} (loop zorgvuldig langs de rand en tel elke zijde mee).
\end{exercise}

\begin{exercise}[$\star\star$]\label{exo:g6:measure:10}
Een zwembad is een \omterm{def:g5:solids:def}{balk} van $25$ m lang, $10$ m breed en $2$ m diep.
\begin{enumerate}
  \item Hoeveel kubieke meter water gaat er in als het vol is?
  \item Hoeveel liter is dat?
  \item Het bad wordt gevuld met $50\,000$ L per uur. Hoe lang duurt het vullen?
\end{enumerate}
\end{exercise}

\begin{exercise}[$\star\star$]\label{exo:g6:measure:11}
Een tuin is een vierkant met zijde $20$ m met daarin een ronde vijver met \omterm{def:g3:shapes:shapes}{radius}
$3$ m. Hoeveel \omterm{def:g6:measure:area}{oppervlakte} gras is er (precieze waarde, en daarna afgerond op de
m$^2$)?
\end{exercise}

\begin{exercise}[$\star\star\star$]\label{exo:g6:measure:12}
Een chocoladereep meet $15$ cm $\times$ $6$ cm $\times$ $1$ cm. De fabrikant
verdubbelt \emph{alle drie} de afmetingen.
\begin{enumerate}
  \item Met hoeveel wordt het \omterm{def:g6:measure:volume}{volume} vermenigvuldigd? Controleer het door beide
        \omterm{def:g6:measure:volume}{volumes} uit te rekenen.
  \item De prijs wordt met $4$ vermenigvuldigd. Is de grote reep voordeliger dan
        de kleine?
\end{enumerate}
\end{exercise}

\section{Opgave: het hek van koningin Dido}

\begin{problem}[{Weekendopgave --- met een vaste lengte hek: welke vorm
omsluit het meeste land? Het vierkant verslaat elke rechthoek, de cirkel
verslaat het vierkant, en een schuurmuur verandert alles}]
\label{pb:g6:measure:1}
De legende zegt dat koningin Dido, toen ze op de kust van Afrika landde,
``zoveel land kreeg als een ossenhuid kan omsluiten'' --- waarop zij de huid in
één immens lange, dunne strook sneed en genoeg grond omsloot om de stad Carthago
te stichten. Haar probleem is nu het jouwe: een boer heeft precies $20$ m hek.
\Cref{exo:g6:measure:8} liet zien dat twee weides met dezelfde \omterm{def:g6:measure:perimeter}{omtrek}
verschillende \omterm{def:g6:measure:area}{oppervlaktes} kunnen hebben; in deze opgave zoek je de \emph{beste}
weide --- en ontdek je dat het antwoord volledig verandert zodra een schuurmuur
een zijde gratis levert.

\medskip
\noindent\textbf{Deel I --- Twintig meter hek.}
\begin{enumerate}
  \item De weide moet een rechthoek zijn die alle $20$ m hek gebruikt. Ga na dat
        een weide van $1 \times 9$ en een van $2 \times 8$ beide meedoen, en
        reken hun \omterm{def:g6:measure:area}{oppervlaktes} uit.
  \item Leg uit waarom de lengte en de breedte van elke geldige weide samen
        $10$ m zijn. Maak daarna de volledige tabel van de weides met hele
        getallen ($1 \times 9$ tot $5 \times 5$) met hun \omterm{def:g6:measure:area}{oppervlaktes}. Welke is
        de beste?
  \item Is het de moeite om decimale zijden te proberen? Reken de \omterm{def:g6:measure:area}{oppervlaktes}
        van de weide $4.5 \times 5.5$ en van de weide $4.9 \times 5.1$ uit en
        vergelijk met het vierkant $5 \times 5$. Wat vermoed je?
  \item Vraag 2 maakte van het hekprobleem een pure getallenvraag: \emph{welk
        paar getallen met \omterm{def:g1:addition:def}{som} $10$ heeft het grootste \omterm{def:g2:mult:def}{product}?} Antwoord met je
        tabel, en formuleer de algemene regel die je opmerkt.
  \item Hier is waarom afwijken van het vierkant altijd verliest, met een schaar
        in plaats van algebra: begin bij de weide $5 \times 5$ en maak er de
        weide $6 \times 4$ van door een strook weg te halen en een andere terug
        te plakken. Welke strook gaat eraf, welke komt erbij, en waarom verliest
        de ruil precies één vierkante meter? Leg uit waarom een volgende stap
        (naar $7 \times 3$) nog meer verliest.
\end{enumerate}

\medskip
\noindent\textbf{Deel II --- De schuur, en de cirkel.}
\begin{enumerate}[resume]
  \item De weide wordt nu tegen een lange schuurmuur gebouwd: de muur vervangt
        één lange zijde, en de $20$ m hek dekt alleen de andere drie zijden.
        Maak een tabel van de weides met hele getallen (breedte $b$, lengte $l$
        langs de muur, $b + l + b = 20$) en hun \omterm{def:g6:measure:area}{oppervlaktes}. Welke weide wint nu
        --- en is dat een vierkant?
  \item Verklaar de winnaar met een spiegel (\cref{ch:g6:symmetry}): spiegel de
        weide in de schuurmuur en bekijk de verdubbelde weide. Hoeveel hek
        gebruikt de verdubbelde weide, welke verdubbelde weide is volgens deel I
        de beste, en wat maakt dat van de oorspronkelijke weide?
  \item Dido bouwde geen rechthoek. Buig de $20$ m hek tot een cirkel: reken met
        $\pi \approx 3.14$ zijn \omterm{def:g3:shapes:shapes}{radius} uit (op de cm) en daarna zijn \omterm{def:g6:measure:area}{oppervlakte}
        (op de m$^2$) (\cref{prop:g6:measure:circle},
        \cref{prop:g6:measure:formulas}). Vergelijk met het vierkant
        $5 \times 5$.
  \item Het omgekeerde probleem: er wordt een weide van precies $36$ m$^2$
        gevraagd, met zo weinig hek als mogelijk. Vergelijk de rechthoeken
        $1 \times 36$, $2 \times 18$, $3 \times 12$, $4 \times 9$ en
        $6 \times 6$: hun omtrekken? Welke vorm is het goedkoopst, en hoe is deze
        vraag het \omterm{def:g6:symmetry:def}{spiegelbeeld} van deel I?
  \item In één zin: waarom zijn blikjes, buizen en watertanks zo vaak
        \emph{rond}?
\end{enumerate}

\medskip
\noindent\textbf{Deel III --- \omterm{def:g6:measure:perimeter}{Omtrek} en \omterm{def:g6:measure:area}{oppervlakte} zijn vreemden.}
\begin{enumerate}[resume]
  \item Zoek een rechthoek met een \omterm{def:g6:measure:perimeter}{omtrek} van meer dan $100$ m en een
        \omterm{def:g6:measure:area}{oppervlakte} van minder dan $1$ m$^2$. (Heel lang en heel dun. Decimalen
        mogen.)
  \item Waar of niet waar: ``van twee figuren heeft die met de grootste \omterm{def:g6:measure:perimeter}{omtrek} de
        grootste \omterm{def:g6:measure:area}{oppervlakte}.'' Geef een tegenvoorbeeld uit deze opgave.
  \item Neem de rechthoek $6 \times 4$ en knip midden in een van de lange zijden
        een vierkante hap van $1 \times 1$ eruit (de hap opent naar buiten, als
        een ontbrekende tand). Reken de \omterm{def:g6:measure:area}{oppervlakte} en de \omterm{def:g6:measure:perimeter}{omtrek} van de gehapte
        figuur uit, en vergelijk beide met het origineel. Wat is er met elk
        gebeurd?
  \item Kaartenmakers kennen een vreemd feit: hoe gedetailleerder de kaart, hoe
        \emph{langer} een kustlijn meet --- elke zoom onthult nieuwe kleine
        happen, en vraag 13 laat zien wat elke hap doet. Leg in één of twee
        zinnen uit waarom ``de lengte van de kust van Bretagne'' een glibberig
        getal is, en ``de \omterm{def:g6:measure:area}{oppervlakte} van Bretagne'' niet.
  \item Het examen van de boer. Zoek met $36$ m hek en de schuurmuur ter
        beschikking de beste rechthoekige weide (gebruik de spiegel van vraag 7),
        geef haar \omterm{def:g6:measure:area}{oppervlakte}, en vergelijk met de beste weide die je met $36$ m
        in het open veld bouwt. Hoeveel levert de schuurmuur de boer op?
\end{enumerate}
\end{problem}
